\documentclass[12pt,a4paper]{report}

%  Encodage et langue
\usepackage[utf8]{inputenc}
\usepackage[T1]{fontenc}
\usepackage[french]{babel}

%  Mise en page
\usepackage[top=2.5cm, bottom=2.5cm, left=3cm, right=2.5cm]{geometry}
\usepackage{setspace}
\onehalfspacing

%  Typographie
\usepackage{lmodern}
\usepackage{microtype}

%  Tableaux
\usepackage{booktabs}
\usepackage{tabularx}
\usepackage{multirow}
\usepackage{array}
\usepackage{longtable}
\usepackage{xcolor}
\usepackage{colortbl}
\usepackage{rotating}

%  Figures
\usepackage{graphicx}
\usepackage{float}

%  Paramètres flottants : chaque figure/tableau occupe sa propre page
\renewcommand{\floatpagefraction}{0.80}
\renewcommand{\topfraction}{0.90}
\renewcommand{\bottomfraction}{0.90}
\setcounter{topnumber}{1}
\setcounter{bottomnumber}{1}
\setcounter{totalnumber}{1}

%  Mathématiques
\usepackage{amsmath}
\usepackage{amssymb}

%  Hyperliens
\usepackage[colorlinks=true, linkcolor=blue, citecolor=blue, urlcolor=blue]{hyperref}

%  Bibliographie
\usepackage[backend=biber, style=numeric, sorting=none]{biblatex}
\addbibresource{references_kalati_rag.bib}

% ============================================================
\begin{document}
% ============================================================

\chapter{Architecture de KALATI-RAG}

% ============================================================
\section{Positionnement du système}
% ============================================================

\subsection*{KALATI et KALATI-RAG : deux couches architecturalement distinctes}

Le portail KALATI constitue l'infrastructure documentaire existante de CAMRAIL. Il centralise
l'ensemble des textes normatifs de la société --- Convention Collective, Instructions Générales
de Sécurité (IGS), procédures de Marche à Vue, règles de Block, circulaires RH, rapports
techniques d'exploitation --- et les rend accessibles au personnel via un mécanisme
d'authentification par matricule et mot de passe, avec une matrice de droits différenciés
par grade et par département. KALATI est cependant un portail de consultation statique~: il
suppose que l'agent dispose d'une connaissance préalable de l'arborescence documentaire,
maîtrise le vocabulaire de recherche propre au corpus ferroviaire camerounais et soit capable
d'interpréter les textes normatifs dans leur contexte réglementaire global. Ces trois prérequis
ne sont pas uniformément distribués dans le personnel CAMRAIL~--- c'est précisément l'asymétrie
que KALATI-RAG est conçu pour corriger.

\begin{figure}[p]
  \centering
  \includegraphics[width=\textwidth]{figures/architecture_general_solution.png}
  \caption{Vue d'ensemble de la solution KALATI-RAG~--- positionnement de la couche IA
  conversationnelle sur l'infrastructure documentaire existante de CAMRAIL.}
  \label{fig:architecture-generale}
\end{figure}

KALATI-RAG est la couche conversationnelle intelligente greffée sur cette infrastructure
existante. Architecturalement, il s'agit d'un système de Retrieval-Augmented Generation (RAG)
hybride combinant recherche sémantique dense et recherche lexicale éparse BM25, dont le corpus
de référence est exclusivement constitué des documents versés dans KALATI. KALATI-RAG ne
remplace pas le portail existant, n'en modifie pas la structure documentaire et n'altère pas
sa matrice de contrôle d'accès. Il s'y connecte comme une surcouche d'interrogation en langage
naturel~: là où KALATI expose des documents à naviguer manuellement, KALATI-RAG expose une
interface de question-réponse où l'agent formule sa requête en français naturel et reçoit une
réponse précise, formulée en langage courant, accompagnée des citations exactes des passages
sources qui la fondent~--- document, section, article, page.

\begin{figure}[p]
  \centering
  \includegraphics[width=\textwidth]{figures/Architecture-globale-KALATI-RAG.png}
  \caption{Architecture globale KALATI-RAG~--- pipeline d'ingestion (Phases~1 et~2) et moteur
  agentique (Phases~3 et~4), greffés sur le portail KALATI sans modification de son
  infrastructure.}
  \label{fig:architecture-globale}
\end{figure}

Cette distinction est architecturalement fondamentale pour une raison épistémologique précise~:
KALATI-RAG ne génère aucune connaissance exogène. Son espace de réponse est borné par le corpus
KALATI. Le modèle de langue ne produit que ce que le retrieveur lui fournit comme contexte~;
si une information est absente du corpus, le système déclenche un message d'indisponibilité
explicite plutôt que de compléter la réponse par des généralités extérieures non vérifiables.
Cette propriété de \textit{grounding} strict est une contrainte de conception délibérée,
directement motivée par le contexte ferroviaire de sécurité dans lequel s'inscrit le système~:
dans un environnement normatif où une réponse incorrecte peut induire une décision opérationnelle
erronée, une réponse absente est préférable à une réponse hallucinée.

\subsection*{Périmètre fonctionnel}

Le périmètre fonctionnel de KALATI-RAG se définit par deux classes de requêtes opérationnelles
et une contrainte de périmètre explicite.

La \textbf{première classe} recouvre les requêtes à passage unique, dans lesquelles l'information
nécessaire à la réponse est localisée dans un seul chunk du corpus. Ce type de requête ---
«~Quel est le délai de préavis applicable à un agent de maîtrise en cas de licenciement selon la
Convention Collective de CAMRAIL~?~» --- mobilise le pipeline de récupération dans sa
configuration nominale~: vectorisation de la question par Qwen3-Embedding-8B, recherche par
similarité cosinus dans l'index Qdrant avec filtre RBAC appliqué en amont, reranking des
passages candidats par Qwen3-Reranker, génération de la réponse par Qwen3-14B en mode
non-thinking. La latence cible pour cette classe est de 4 à 5 secondes, ce qui correspond à
l'essentiel des interactions quotidiennes des agents CAMRAIL.

La \textbf{deuxième classe} recouvre les requêtes multi-hop, dans lesquelles la construction
d'une réponse complète et correcte nécessite l'agrégation de preuves distribuées sur plusieurs
passages de sources documentaires hétérogènes. La littérature récente sur le RAG avancé
identifie l'agrégation de preuves multi-sources comme le défi central des systèmes RAG en
production, les LLMs généralistes présentant des défaillances systématiques de
\textit{grounding} lorsque la réponse requiert de croiser des sources hétérogènes dans des
documents longs~\autocite{Gao2025}. KALATI-RAG adresse cette classe par l'activation automatique
du mode thinking de Qwen3-30B-A3B, qui implémente une chaîne de pensée interne décomposant
la question en sous-questions avant d'effectuer plusieurs passes de récupération itérative.
La latence cible pour cette classe est de 14 à 15 secondes.

La \textbf{contrainte de périmètre} est structurellement aussi importante que les deux classes
de capacités~: KALATI-RAG opère en monde fermé sur le corpus KALATI. Toute question dont la
réponse est absente du corpus déclenche un refus calibré explicite, activé automatiquement
par le module de validation anti-hallucination dès que le score de faithfulness calculé en
temps réel est inférieur au seuil de production fixé à 0,90. Ce mécanisme est décrit en détail
à la section~\ref{sec:validation}.

\subsection*{Périmètre technique~: intranet fermé, zéro flux sortant, zéro coût récurrent}

Le périmètre technique de KALATI-RAG repose sur trois propriétés non négociables qui ont
conditionné l'intégralité des choix architecturaux décrits dans les sections suivantes.

\textbf{Première propriété --- déploiement en intranet fermé.} L'intégralité de la pile
d'inférence --- modèle de langue Qwen3, moteur d'embedding Qwen3-Embedding-8B, reranker
Qwen3-Reranker, base vectorielle Qdrant, moteur BM25, serveur d'orchestration LangChain,
interface web --- est déployée sur les serveurs internes de CAMRAIL et exposée exclusivement
sur le réseau intranet de l'entreprise. Aucun composant du pipeline ne dispose d'une route
réseau vers l'extérieur. Cette propriété est garantie par le périmètre réseau lui-même,
ce qui la rend physiquement non contournable~: la confidentialité des données normatives,
des questions posées par les agents et des journaux d'interaction est ainsi garantie par
architecture, et non par stipulation contractuelle.

\textbf{Deuxième propriété --- zéro flux sortant de données.} Aucune requête d'agent, aucun
extrait du corpus KALATI, aucun vecteur d'embedding et aucun journal d'interaction ne transite
vers un service externe. Cette contrainte exclut par construction toutes les architectures
reposant sur des appels à des API distantes telles qu'OpenAI, Anthropic, Cohere ou Mistral API,
dont les conditions d'utilisation autorisent le traitement des données transmises sur des
infrastructures tierces. Elle est le fondement de la souveraineté de CAMRAIL sur sa propre
connaissance normative et réglementaire.

\textbf{Troisième propriété --- zéro coût récurrent de tokens.} La pile technologique retenue
est intégralement composée de solutions open source distribuées sous licence Apache~2.0,
déployées en auto-hébergement sans quota d'utilisation ni frais de licence. Qwen3-14B,
Qwen3-Embedding-8B, Qwen3-Reranker, Qdrant, Ollama et LangChain sont tous librement
utilisables dans un contexte institutionnel sans restriction commerciale. Une fois le système
mis en production, son coût opérationnel marginal est nul du côté logiciel~: il se réduit
aux coûts d'infrastructure matérielle et de maintenance humaine, deux postes que CAMRAIL
dimensionne et maîtrise en interne.

% ============================================================
\section{Phase~1 --- Ingestion et traitement du corpus KALATI}
% ============================================================

\subsection{Classification des types documentaires}

La première opération du pipeline KALATI-RAG consiste à analyser la nature de chaque document
entrant afin de lui appliquer le traitement d'extraction adapté. Le corpus KALATI présente une
hétérogénéité structurelle qui interdit toute approche d'ingestion uniforme~: une circulaire RH
produite sous traitement de texte n'a pas les mêmes propriétés d'extraction qu'un rapport
technique issu d'une numérisation d'archives papier. Une analyse préalable du fonds documentaire
de CAMRAIL permet d'identifier cinq classes de documents aux propriétés radicalement différentes,
auxquelles correspondent cinq traitements distincts dans le pipeline.

\begin{figure}[p]
  \centering
  \includegraphics[width=\textwidth]{figures/KALATI-Document-Processing-Flowchart.png}
  \caption{Arbre de décision du routage documentaire KALATI-RAG~--- du document KALATI au
  traitement approprié selon ses propriétés structurelles~: extraction directe, OCR conditionnel
  ou extraction structurée des tableaux.}
  \label{fig:flowchart}
\end{figure}

\begin{sidewaystable}[p]
\centering
\caption{Classification des types documentaires du corpus KALATI et traitements associés}
\label{tab:classes-docs}
\small
\renewcommand{\arraystretch}{1.2}
\begin{tabularx}{\textheight}{@{} p{4.5cm} X p{9cm} @{}}
\toprule
\textbf{Classe} & \textbf{Description} & \textbf{Traitement} \\
\midrule
PDF natif à texte extractible &
Produits par outils bureautiques~; couche textuelle encodée dans le fichier PDF~; inclut les
textes normatifs récents, circulaires RH, procédures numériques &
Extraction directe par PyMuPDF (\texttt{fitz})~--- aucune inférence neuronale, fidélité de
transcription 100~\% sur les caractères encodés \\
\addlinespace
PDF image (archives numérisées) &
Scans ou photographies de documents papier sans couche textuelle~; archives ferroviaires
anciennes avec dégradations géométriques et photométriques &
Chaîne OCR~: PaddleOCR-VL~1.6 (moteur principal, OmniDocBench~96,33~\%~\autocite{Shi2024}) \\
\addlinespace
Documents tamponnés &
Circulaires, notes de service et décisions comportant des cachets officiels apposés manuellement,
superposés au texte imprimé &
Détection de formes circulaires par transformée de Hough, puis routage vers GLM-OCR (précision
tampons~90,5~\%~\autocite{ZhipuAI2025}) \\
\addlinespace
Documents contenant des tableaux structurés &
Grilles tarifaires, matrices de droits par grade, tableaux de procédures de maintenance dont
la structure relationnelle doit être préservée lors de l'indexation &
Table Transformer (TATR~\autocite{Smock2022})~--- extraction cellule par cellule avec
reconstruction JSON préservant les correspondances en-tête/valeur \\
\addlinespace
Documents contenant des images embarquées &
Schémas de voie ferrée, diagrammes de signalisation, plans de maintenance dont le contenu est
purement graphique &
PaddleOCR-VL en mode vision-langage~--- génération de description textuelle structurée et
indexable du contenu visuel \\
\bottomrule
\end{tabularx}
\end{sidewaystable}

\subsection{Extraction directe des PDF natifs}

Pour les documents de la première classe, le pipeline s'affranchit de toute inférence neuronale
et recourt à une extraction programmatique directe de la couche textuelle encodée dans le fichier
PDF. Cette extraction est réalisée par la bibliothèque PyMuPDF (\texttt{fitz}), qui permet
l'accès à la représentation interne du document PDF selon la spécification ISO~32000 et restitue
le texte avec sa structure de blocs, de lignes et de caractères, ainsi que les métadonnées de
position spatiale de chaque élément textuel sur la page.

Cette approche présente trois avantages déterminants sur le traitement OCR pour cette classe~: une
fidélité de transcription de 100~\% sur les caractères encodés, une vitesse d'extraction de
plusieurs ordres de grandeur supérieure à celle de l'inférence neuronale, et l'absence de tout
risque de substitution de caractères. Pour un corpus normatif où la précision lexicale est une
contrainte de sécurité --- un article de la Convention Collective ne peut pas être transcrit avec
une erreur sur un chiffre ou un terme technique ---, cette classe de documents est traitée sans
passer par la chaîne OCR. Un test de détection de la couche textuelle est appliqué
systématiquement à chaque document entrant~: si le ratio caractères extractibles sur surface de
page dépasse un seuil calibré lors de la Phase~0, le document est classé comme PDF natif~; dans
le cas contraire, il est routé vers la chaîne OCR décrite à la section suivante.

\subsection{Sélection de la chaîne OCR par benchmark}

\subsubsection*{Référentiel d'évaluation~: OmniDocBench (CVPR~2025)}

La sélection des moteurs OCR repose sur le benchmark OmniDocBench, publié par \textcite{Shi2024}
à la conférence CVPR~2025, qui constitue à ce jour la référence la plus rigoureuse et la plus
exhaustive pour l'évaluation comparative des moteurs de traitement documentaire en conditions
réelles. OmniDocBench se distingue des benchmarks antérieurs par trois propriétés
méthodologiques déterminantes pour le contexte CAMRAIL~: il évalue les moteurs sur des documents
hétérogènes couvrant dix-neuf types documentaires distincts~; il intègre des conditions de
dégradation documentaire réalistes (distorsion géométrique, variation d'illumination, bruit de
fond) qui correspondent précisément aux conditions des archives physiques de CAMRAIL~; et il
mesure la performance de bout en bout sur la qualité de la transcription textuelle, et non
seulement sur la détection de régions d'intérêt.

\begin{figure}[p]
  \centering
  \includegraphics[width=\textwidth]{figures/Pipeline-OCR.png}
  \caption{Pipeline OCR conditionnel KALATI-RAG~--- du PDF image au texte normalisé prêt à
  segmenter, avec routage conditionnel sur la détection de tampons administratifs par transformée
  de Hough.}
  \label{fig:pipeline-ocr}
\end{figure}

\subsubsection*{Panorama des candidats évalués et sélection}

\begin{table}[p]
\centering
\caption{Comparaison des moteurs OCR candidats sur OmniDocBench~\autocite{Shi2024} (CVPR~2025)}
\label{tab:ocr-benchmark}
\renewcommand{\arraystretch}{1.4}
\begin{tabularx}{\linewidth}{@{} l c X l @{}}
\toprule
\textbf{Moteur} & \textbf{Score} & \textbf{Caractéristiques clés} & \textbf{Décision} \\
\midrule
PaddleOCR-VL~1.6 \newline (Baidu, Apache~2.0) & 96,33~\% &
Architecture vision-langage~; robustesse aux dégradations~; 109~langues~; compatible CPU &
\textbf{Retenu} \\
\addlinespace
GLM-OCR \newline (Zhipu AI, MIT) & 94,62~\% &
Précision 90,5~\% sur tampons~; 0,9~Md paramètres~; 2--4~Go VRAM &
\textbf{Retenu} (tampons) \\
\addlinespace
docTR \newline (Mindee, Apache~2.0) & 78,1~\% &
Pipeline DBNet + CRNN sans modélisation multimodale &
Éliminé \\
\addlinespace
Tesseract~5 \newline (Google, Apache~2.0) & 72,4~\% &
LSTM sur corpus synthétiques~; $<$60~\% sur zones tamponnées &
Éliminé \\
\addlinespace
DeepSeek-OCR2 & n/a &
Performant mais licence incompatible avec intranet institutionnel fermé &
Éliminé \\
\bottomrule
\end{tabularx}
\end{table}

\subsubsection*{Détection des tampons et nœud de décision}

La détection préalable de la présence d'un tampon administratif est réalisée par la transformée
de Hough appliquée aux contours extraits de l'image du document. Cette méthode de traitement
d'image classique, sans inférence neuronale, détecte les formes circulaires et ovales
caractéristiques des cachets officiels à un coût computationnel de l'ordre de quelques
millisecondes par page. Si la transformée de Hough détecte au moins une forme circulaire
d'une superficie supérieure au seuil calibré en Phase~0, le document est routé vers GLM-OCR~;
dans le cas contraire, il est traité par PaddleOCR-VL~1.6. Les deux moteurs convergent ensuite
vers la même étape de post-correction.

\subsection{Extraction des tableaux comme données structurées}

Le corpus KALATI contient un nombre significatif de documents intégrant des tableaux structurés
--- grilles tarifaires, matrices de droits par grade, tableaux de procédures de maintenance
ferroviaire. La transcription de ces tableaux par un moteur OCR généraliste produit une séquence
de tokens linéaire qui détruit la structure relationnelle~: les correspondances entre en-têtes
de colonnes et valeurs de cellules sont perdues, rendant les données inexploitables pour une
requête de type «~Quel est le délai de préavis d'un agent de maîtrise~?~» lorsque la réponse
figure dans un tableau récapitulatif. Afin de préserver cette structure, KALATI-RAG intègre le
\textbf{Table Transformer (TATR)}, publié par \textcite{Smock2022} à CVPR~2022, comme module
de détection et d'extraction structurée.

\begin{figure}[p]
  \centering
  \includegraphics[width=\textwidth]{figures/extraction_table_par_ocr.png}
  \caption{Pipeline d'extraction structurée des tableaux par Table Transformer (TATR)~---
  détection des régions, segmentation cellule par cellule, reconstruction JSON et indexation dans
  Qdrant avec métadonnées \texttt{row\_header} et \texttt{col\_header}.}
  \label{fig:tatr}
\end{figure}

TATR est un modèle de détection d'objets basé sur DETR (\textit{DEtection TRansformer}) entraîné
sur PubTables-1M pour détecter les régions de tableaux, les lignes, les colonnes et les cellules
individuelles avec leurs coordonnées spatiales. Le pipeline d'extraction structurée opère en
quatre étapes~: (1)~TATR identifie les bounding boxes des tableaux présents dans chaque page~;
(2)~segmentation cellule par cellule avec attribution des coordonnées (ligne, colonne)~;
(3)~reconstruction de la structure relationnelle en JSON avec les en-têtes de colonnes comme
clés~; (4)~indexation dans Qdrant comme chunk de type \texttt{cell}, avec les métadonnées
\texttt{row\_header} et \texttt{col\_header} permettant de reconstruire la réponse structurée
en génération. Ce traitement garantit que la requête «~Quel est le délai de préavis d'un agent
de maîtrise~?~» retrouve non seulement le paragraphe textuel de l'article concerné, mais aussi
la ligne correspondante du tableau récapitulatif, les deux sources se renforçant mutuellement.

\subsection{Traitement des images embarquées dans le corpus}

Les documents du corpus KALATI contiennent des images techniques embarquées~--- schémas de voie
ferrée, diagrammes de signalisation, plans de maintenance, organigrammes des procédures
d'exploitation. Ces images sont absentes du texte transcrit par la chaîne OCR standard, qui ne
produit aucun token à partir d'un contenu purement graphique. Pour remédier à cette lacune, le
traitement des images embarquées repose sur les capacités vision-langage de PaddleOCR-VL~1.6 en
mode multimodal, qui génère une description textuelle structurée de chaque schéma ou diagramme
technique. La description générée est indexée dans Qdrant comme chunk de type \texttt{image\_caption},
avec la métadonnée \texttt{image\_type} (schéma, diagramme, organigramme, photographie)
utilisable comme filtre optionnel lors de la récupération.

\subsection{Post-correction par le lexique KALATI}

La transcription OCR, même à des scores OmniDocBench supérieurs à 94~\%, introduit des
substitutions de caractères statistiquement prévisibles sur les termes à faible fréquence dans
les corpus d'entraînement des moteurs. Le vocabulaire ferroviaire camerounais~--- sigles
d'exploitation (IGS, PTU, PCC, MAV pour Marche à Vue), abréviations de direction, noms de gares
et de tronçons du réseau CAMRAIL~--- appartient précisément à cette catégorie de termes absents
des corpus généraux. La Phase~0 produit un lexique ferroviaire CAMRAIL constitué par extraction
semi-automatique des termes techniques présents dans un échantillon de référence du corpus. Ce
lexique est utilisé en Phase~1 comme dictionnaire de correction post-OCR~: chaque terme du texte
transcrit est comparé aux entrées du lexique par distance d'édition de Levenshtein normalisée, et
les substitutions dont le score de similarité dépasse un seuil calibré sont corrigées
automatiquement. Cette étape garantit notamment que «~Marche à Vue~» ne devient jamais
«~Marche à Voie~» dans le corpus indexé.

\subsection{Segmentation en chunks}

La segmentation des textes post-corrigés en unités indexables constitue l'une des décisions
architecturales les plus déterminantes pour la qualité de la récupération en Phase~3. La
littérature sur les systèmes RAG en production établit un compromis fondamental entre deux
propriétés antagonistes~: la précision de la localisation de l'information, favorisée par des
chunks courts, et la qualité du contexte fourni au modèle de langue lors de la génération,
favorisée par des chunks longs.

\begin{figure}[p]
  \centering
  \includegraphics[width=\textwidth]{figures/chuncking_indexing.png}
  \caption{Stratégie de double granularité (\textit{parent-child chunking})~--- les chunks courts
  (128 tokens, en vert) servent à la récupération précise~; les chunks longs correspondants
  (512 tokens, en bleu) sont transmis au LLM comme contexte de génération.}
  \label{fig:chunking}
\end{figure}

KALATI-RAG résout ce compromis par une stratégie de double granularité~: chaque segment de texte
post-corrigé est simultanément découpé en chunks courts de 128 tokens (unité de recherche) et en
chunks longs de 512 tokens (unité de génération), avec un chevauchement de 20 tokens entre chunks
adjacents pour préserver la continuité sémantique aux frontières. Cette architecture, désignée
dans la littérature sous le terme de \textit{parent-child chunking}~\autocite{Chen2023}, dissocie
la précision de la recherche (chunks courts) de la richesse du contexte de génération (chunks
longs correspondants). Pour les documents à structure d'articles bien définie~--- Convention
Collective, IGS~---, une stratégie complémentaire de \textbf{semantic chunking}~\autocite{Chen2023}
est appliquée~: plutôt que de découper le texte en fenêtres de tokens fixes, les frontières de
segmentation sont déterminées par les marqueurs structurels du document (numéros d'articles,
tirets de liste, paragraphes), produisant des chunks dont la complétude sémantique est garantie
par la structure documentaire elle-même.

\subsection{Enrichissement des métadonnées héritées de la taxonomie KALATI}

Chaque chunk produit à l'issue de la segmentation est enrichi d'un ensemble de métadonnées
structurées héritées de la taxonomie documentaire du portail KALATI. Ces métadonnées sont stockées
comme payload aux côtés du vecteur d'embedding dans Qdrant et jouent un double rôle~: elles
permettent la traçabilité des sources dans les réponses générées, et elles constituent le support
du filtre RBAC appliqué lors de chaque requête d'agent.

\begin{table}[p]
\centering
\caption{Métadonnées associées à chaque chunk dans Qdrant}
\label{tab:metadata}
\renewcommand{\arraystretch}{1.4}
\begin{tabularx}{\linewidth}{@{} p{3.5cm} p{2.5cm} X @{}}
\toprule
\textbf{Champ} & \textbf{Type} & \textbf{Rôle} \\
\midrule
\texttt{doc\_id} & Chaîne & Identifiant unique du document source dans KALATI \\
\texttt{titre} & Chaîne & Titre du document source \\
\texttt{categorie} & Énumération & IGS, Convention Collective, circulaire RH, procédure technique \\
\texttt{departement} & Chaîne & Direction émettrice selon l'organigramme CAMRAIL \\
\texttt{grade\_min} & Entier & Niveau d'habilitation minimal~--- critère de filtrage RBAC \\
\texttt{page} & Entier & Numéro de page dans le document source \\
\texttt{section} & Chaîne & Référence de section ou d'article \\
\texttt{chunk\_type} & Énumération & \texttt{short}, \texttt{long}, \texttt{cell},
\texttt{image\_caption}, \texttt{graph\_summary} \\
\texttt{status} & Énumération & \texttt{active} ou \texttt{superseded}~--- filtre systématique \\
\texttt{superseded\_by} & Chaîne & Identifiant \texttt{doc\_id} du document remplaçant \\
\bottomrule
\end{tabularx}
\end{table}

Le champ \texttt{grade\_min} est le champ central du mécanisme RBAC~: lors de chaque requête en
Phase~3, le filtre Qdrant construit à partir des claims du token d'authentification de l'agent
impose la condition \texttt{grade\_min}~$\leq$~\texttt{grade\_agent}, garantissant que seuls les
passages auxquels l'agent est habilité sont retournés par le moteur de recherche.

\subsection{Gestion des versions documentaires et mises à jour incrémentales}

Le corpus KALATI est un corpus vivant~: les textes normatifs de CAMRAIL font l'objet de révisions
régulières~--- nouvelles circulaires RH annulant des circulaires antérieures, mises à jour des
IGS, avenants à la Convention Collective. Le schéma d'identifiant documentaire retenu intègre
la version comme composante structurelle~: \texttt{doc\_id = "\{type\_doc\}\_\{code\_doc\}\_\{version\}\_\{date\_effet\}"}
par exemple \texttt{IGS\_CAMRAIL04\_v3\_20260101} pour la troisième révision de l'IGS CAMRAIL04.

Lors d'une mise à jour documentaire, les chunks de l'ancienne version sont marqués
\texttt{superseded} et leur champ \texttt{superseded\_by} est renseigné avec l'identifiant de
la nouvelle version~; les chunks de la nouvelle version sont insérés avec le statut
\texttt{active}. Le filtre Qdrant appliqué lors de chaque requête impose systématiquement la
condition \texttt{status = active}, garantissant qu'un agent ne peut pas recevoir une réponse
fondée sur une version obsolète d'un texte normatif. Les versions remplacées restent présentes
dans la collection avec le statut \texttt{superseded} pour permettre les requêtes historiques
formulées par un administrateur~--- par exemple, pour vérifier la disposition applicable à la
date d'un incident passé.

% ============================================================
\section{Phase~2 --- Vectorisation et indexation}
% ============================================================

\subsection{Sélection du modèle d'embedding}

\subsubsection*{Référentiel d'évaluation~: MTEB Multilingual Leaderboard 2026}

La sélection du modèle d'embedding repose sur le Massive Text Embedding Benchmark (MTEB),
publié par \textcite{Muennighoff2023} et maintenu comme référentiel d'évaluation standard de la
communauté depuis 2023. Le MTEB Multilingual Leaderboard 2026 constitue à ce jour l'évaluation
comparative la plus exhaustive des modèles d'embedding multilingues~: il agrège les performances
sur 56~jeux de données couvrant huit tâches distinctes dans plus de 112~langues. Ce référentiel
est particulièrement adapté au contexte KALATI pour deux raisons~: la tâche de retrieval
bénéficie d'une évaluation dédiée sur des corpus multilingues incluant le français~; et le MTEB
évalue les modèles en conditions zero-shot sur des domaines non vus lors de l'entraînement~---
condition précise d'utilisation dans KALATI-RAG, le corpus ferroviaire camerounais étant absent
des données d'entraînement de tout modèle d'embedding disponible.

\begin{sidewaystable}[p]
\centering
\caption{Comparaison des modèles d'embedding candidats (MTEB Multilingual Leaderboard 2026,
\autocite{Muennighoff2023})}
\label{tab:embedding-benchmark}
\renewcommand{\arraystretch}{1.4}
\begin{tabularx}{\textheight}{@{} l c c X l @{}}
\toprule
\textbf{Modèle} & \textbf{Dim.} & \textbf{Score MTEB} & \textbf{Limitation principale}
& \textbf{Décision} \\
\midrule
Qwen3-Embedding-8B \newline (Alibaba, Apache~2.0) & 3~072 & Meilleur open source 2026 &
Aucune~: meilleure performance multilangue, écosystème Qwen3 aligné avec LLM et Reranker nativement &
\textbf{Retenu} \\
\addlinespace
bge-m3 (BAAI, MIT) & --- & 68,2~\% &
Pas d'écosystème LLM/reranker aligné~: discontinuité sémantique inter-modèles documentée &
Éliminé \\
\addlinespace
multilingual-e5-large \newline (Microsoft, MIT) & 1~024 & 61,5~\% &
Reranker natif non aligné~; dimensionnalité insuffisante pour corpus normatif spécialisé &
Éliminé \\
\addlinespace
paraphrase-multilingual-mpnet-base-v2 \newline (Apache~2.0) & 768 & 53,4~\% &
Performances retrieval insuffisantes sur domaine spécialisé en français &
Éliminé \\
\bottomrule
\end{tabularx}
\end{sidewaystable}

Qwen3-Embedding-8B est retenu pour deux propriétés techniques déterminantes. Premièrement,
ses performances sur la tâche de retrieval multilingue en français sont parmi les meilleures
documentées dans la catégorie open source. Deuxièmement~--- et c'est la propriété la plus
importante architecturalement~--- il partage le même espace de représentation sémantique que
les modèles de langue Qwen3-14B et Qwen3-30B-A3B, éliminant la discontinuité sémantique
inter-modèles décrite à la section suivante.

\subsection{Cohérence écosystémique LLM~-- Embedding~-- Reranker}
\label{sec:cohesion}

La propriété architecturale la plus structurante du pipeline KALATI-RAG est la cohérence
écosystémique de bout en bout entre ses trois composants de traitement sémantique~: le modèle
de langue Qwen3-14B (ou Qwen3-30B-A3B), le modèle d'embedding Qwen3-Embedding-8B et le
reranker Qwen3-Reranker. Ces trois composants appartiennent à la même famille de modèles,
entraînés sur des corpus alignés et partageant le même espace de représentation sémantique.

\begin{figure}[p]
  \centering
  \includegraphics[width=\textwidth]{figures/coherence_ecosystemique.png}
  \caption{Cohérence écosystémique Qwen3~--- espace sémantique partagé entre LLM, embedding
  et reranker (droite), par opposition à un assemblage hétérogène de modèles indépendants dont
  les espaces de représentation sont désalignés (gauche).}
  \label{fig:coherence}
\end{figure}

Dans un pipeline RAG hétérogène~--- par exemple, un LLM Mistral couplé à un embedding bge-m3 et
un reranker cross-encoder indépendant~---, chaque composant a été entraîné dans un espace
sémantique distinct. Cette discontinuité entre retrieval et génération est une source documentée
de dégradation de la faithfulness dans les systèmes RAG en production~: le modèle de langue
reçoit des passages dont la représentation sémantique n'est pas optimisée pour son propre espace
d'attention. Dans l'architecture KALATI-RAG, en revanche, le vecteur produit par
Qwen3-Embedding-8B pour un chunk du corpus et le vecteur produit pour une requête d'agent sont
dans le même espace sémantique que celui dans lequel Qwen3-14B a été entraîné à interpréter les
passages fournis comme contexte. Le score de similarité cosinus calculé par Qdrant est donc
directement interprétable par le modèle de langue comme une mesure de pertinence dans son propre
espace représentationnel. Cette cohérence de bout en bout maximise la faithfulness de la
génération sans nécessiter de couche d'adaptation supplémentaire entre les composants.

\subsection{Sélection et benchmark du moteur vectoriel Qdrant}

\subsubsection*{Enjeux de la sélection pour le contexte KALATI-RAG}

Le choix du moteur de base vectorielle est une décision d'infrastructure aussi structurante que
le choix du modèle de langue. Dans le contexte spécifique de KALATI-RAG, quatre critères non
négociables conditionnent cette sélection~: (1)~le filtrage RBAC natif au niveau du moteur, sans
surcoût de latence, afin de garantir que chaque requête n'accède qu'aux passages auxquels l'agent
est habilité~; (2)~la prise en charge native de la recherche hybride dense + BM25, condition du
pipeline de fusion RRF décrit à la section suivante~; (3)~les performances sous filtrage en
conditions de corpus à grande échelle, critère distinctif documenté par les benchmarks
comparatifs~; (4)~le déploiement en auto-hébergement sous licence open source compatible avec
l'intranet fermé de CAMRAIL.

\subsubsection*{Référentiel d'évaluation~: VectorDBBench (2024)}

La comparaison des moteurs vectoriels candidats repose sur \textbf{VectorDBBench}~\autocite{Pan2024},
le benchmark open source de référence pour l'évaluation comparative des systèmes de gestion de
bases de données vectorielles. VectorDBBench mesure trois dimensions clés~: le recall@10 (taux
de rappel sur les dix passages les plus pertinents), les requêtes par seconde (QPS) en condition
de filtrage de métadonnées~--- cas représentatif du filtrage RBAC de KALATI-RAG~---, et la
consommation mémoire à l'indexation. Les mesures sont réalisées sur une collection de 1~million
de vecteurs denses, correspondant à l'ordre de grandeur estimé du corpus KALATI après ingestion
complète.

\begin{sidewaystable}[p]
\centering
\caption{Benchmark comparatif des moteurs vectoriels candidats (VectorDBBench~2024,
\autocite{Pan2024})~--- mesures sur collection de 1~million de vecteurs}
\label{tab:qdrant-benchmark}
\renewcommand{\arraystretch}{1.5}
\begin{tabularx}{\textheight}{@{} l c c c c c c @{}}
\toprule
\textbf{Critère} & \textbf{Qdrant} & \textbf{Weaviate} & \textbf{Milvus} & \textbf{Chroma}
& \textbf{pgvector} & \textbf{FAISS} \\
\midrule
Recall@10 (filtre actif) & \textbf{0,989} & 0,974 & 0,982 & n/a & 0,923 & 0,978 \\
\addlinespace
QPS sous filtrage RBAC & \textbf{850} & 320 & 610 & n/a & 145 & --- \\
\addlinespace
Mémoire / 1M vect. (3~072~dims) & \textbf{13~Go} & 18~Go & 16~Go & 22~Go & 30~Go & 12~Go \\
\addlinespace
Filtrage payload natif & \textbf{\checkmark} & \checkmark & \checkmark & \texttimes
& \checkmark & \texttimes \\
\addlinespace
Recherche hybride dense + BM25 native & \textbf{\checkmark} & \checkmark & partielle
& \texttimes & \texttimes & \texttimes \\
\addlinespace
RBAC appliqué au niveau du moteur & \textbf{\checkmark} & \texttimes & \texttimes & \texttimes
& \texttimes & \texttimes \\
\addlinespace
Déploiement intranet auto-hébergé & \textbf{\checkmark} & \checkmark & \checkmark & \checkmark
& \checkmark & \checkmark \\
\addlinespace
Licence & \textbf{Apache~2.0} & BSD-3 & Apache~2.0 & Apache~2.0 & PostgreSQL & MIT \\
\addlinespace
Intégration LangChain native & \textbf{\checkmark} & \checkmark & \checkmark & \checkmark
& \checkmark & \checkmark \\
\addlinespace
\textbf{Décision} & \textbf{Retenu} & Éliminé & Éliminé & Éliminé & Éliminé & Éliminé \\
\bottomrule
\end{tabularx}
\end{sidewaystable}

\subsubsection*{Justification de la sélection de Qdrant}

L'analyse comparative révèle que Qdrant se distingue des autres candidats sur le critère le plus
structurant pour KALATI-RAG~: le \textbf{filtrage RBAC natif au niveau du moteur}. Aucun des
autres systèmes évalués ne permet d'appliquer des contraintes d'accès directement dans le moteur
de recherche vectorielle~; ils délèguent ce filtrage à la couche applicative, créant une fenêtre
de vulnérabilité entre la récupération et le filtrage. Chez Qdrant, en revanche, le filtre RBAC
est une clause \texttt{must} jointe à la requête HNSW~: il s'applique \emph{avant} toute
récupération de passage, rendant architecturalement impossible l'accès d'un agent à un passage
pour lequel son grade ou son département ne lui confère pas d'habilitation.

Par ailleurs, Qdrant atteint le recall@10 le plus élevé en condition de filtrage actif (0,989)
tout en maintenant le QPS le plus élevé (850~requêtes/seconde), démontrant que ses optimisations
d'index HNSW préservent leurs performances même lorsqu'un filtre de métadonnées réduit
significativement l'espace de recherche~--- limitation documentée chez pgvector, dont le recall@10
chute à 0,923 sous filtrage en raison de l'absence d'index vectoriel nativement optimisé pour
les requêtes filtrées. Enfin, la prise en charge native de la recherche hybride dense + BM25 et
son implémentation native de la RRF élimine tout développement custom d'une couche de fusion,
réduisant ainsi la surface de maintenance du système.

\subsubsection*{Architecture interne de l'index Qdrant}

\begin{figure}[p]
  \centering
  \includegraphics[width=\textwidth]{figures/architechture_interne_Qdran_and_rbac_filter.png}
  \caption{Structure interne de l'index Qdrant~--- graphe HNSW multi-couches pour la recherche
  approximative rapide (haut) et payload RBAC attaché à chaque vecteur pour le filtrage d'accès
  avant récupération (bas).}
  \label{fig:qdrant-rbac}
\end{figure}

Les vecteurs denses sont indexés en 3~072 dimensions, correspondant à la dimensionnalité native
de Qwen3-Embedding-8B. L'index est construit selon l'algorithme HNSW
(\textit{Hierarchical Navigable Small World}), qui offre un compromis optimal entre vitesse de
recherche par similarité approximative et précision du rappel~: complexité de recherche en
$O(\log n)$, configurable via les paramètres $m$ (nombre de connexions par nœud) et
$ef\_construction$ (taille de la file de candidats lors de la construction), tous deux calibrés
lors de la Phase~0 en fonction du volume documentaire effectif du corpus. La métrique de
similarité retenue est la similarité cosinus, cohérente avec le mode d'entraînement de
Qwen3-Embedding-8B. Chaque vecteur indexé est accompagné d'un payload JSON structuré stockant
les métadonnées documentaires décrites à la section~2.7~; ce payload est indexé par Qdrant
indépendamment du vecteur, permettant d'appliquer des filtres sur les champs de métadonnées
simultanément à la recherche vectorielle sans surcoût de latence significatif. Le chargement
est réalisé par insertion en batch asynchrone avec vérification d'idempotence sur \texttt{doc\_id},
garantissant la cohérence lors des mises à jour incrémentales sans nécessiter de réindexation
complète.

\subsection{Fusion hybride BM25 -- dense~: Reciprocal Rank Fusion}

La recherche hybride de KALATI-RAG combine deux signaux de pertinence de nature fondamentalement
différente~: la recherche dense par similarité cosinus dans l'index HNSW de Qdrant, et la
recherche lexicale éparse par BM25 sur le corpus tokenisé. Ces deux signaux produisent chacun
une liste ordonnée de passages candidats dont les scores ne sont pas directement comparables~---
le score cosinus est borné dans $[-1, 1]$, le score BM25 est un score TF-IDF non borné calibré
sur la distribution du corpus. La fusion de ces deux listes est réalisée par
\textbf{Reciprocal Rank Fusion (RRF)}, proposée par \textcite{Cormack2009}~:

\begin{equation}
\text{RRF}(d) = \sum_{i \in \{\text{dense},\, \text{BM25}\}} \frac{1}{k + \text{rank}_i(d)}
\label{eq:rrf}
\end{equation}

où $\text{rank}_i(d)$ est le rang de $d$ dans la liste $i$ (rang~1 = score le plus élevé) et
$k = 60$ est la constante de lissage standardisée dans la littérature. La RRF présente trois
propriétés déterminantes pour KALATI-RAG~: (1)~elle opère sur les rangs et non sur les scores
bruts, la rendant robuste aux distributions de scores très différentes entre BM25 et cosinus~;
(2)~elle assure la complémentarité lexicale-sémantique~--- un passage contenant les termes exacts
de la requête (BM25 élevé) et un passage sémantiquement proche sans correspondance lexicale exacte
(cosinus élevé) contribuent tous deux de manière équilibrée au score final~; (3)~elle est
implémentée nativement dans l'API de recherche hybride de Qdrant, sans surcoût d'implémentation.

\begin{figure}[p]
  \centering
  \includegraphics[width=\textwidth]{figures/hybrid_fontion_bm25_x_RRF.png}
  \caption{Mécanisme de fusion hybride BM25 + dense par Reciprocal Rank Fusion~--- les deux
  listes ordonnées issues des moteurs lexical et sémantique convergent dans un score de rang
  agrégé unique, implémenté nativement dans Qdrant.}
  \label{fig:rrf}
\end{figure}

Cette complémentarité est particulièrement pertinente pour le corpus KALATI~: une requête en
langage naturel courant («~agent qui commence son travail~») doit retrouver des passages
contenant le terme normatif précis («~agent en période d'essai~»), ce que le signal BM25 seul
ne peut pas faire. Inversement, une requête utilisant le sigle exact («~MAV~») doit retrouver
tous les passages relatifs à la Marche à Vue, y compris ceux qui emploient l'expression
développée~--- rôle assuré par le signal dense.

% ============================================================
\section{Phase~3 --- Moteur RAG agentique}
% ============================================================

\subsection{Sélection du modèle de langue~: démarche comparative}

\subsubsection*{Formalisation des contraintes CAMRAIL comme grille de sélection}

Avant toute comparaison de benchmarks, cinq contraintes non négociables, issues de l'analyse
des besoins conduite en Phase~0, constituent la grille de sélection appliquée à l'ensemble des
candidats évalués. Ces contraintes sont formalisées ci-dessous~:

\begin{table}[p]
\centering
\caption{Contraintes CAMRAIL pour la sélection du modèle de langue}
\label{tab:contraintes}
\renewcommand{\arraystretch}{1.4}
\begin{tabularx}{\linewidth}{@{} l l X @{}}
\toprule
\textbf{Code} & \textbf{Contrainte} & \textbf{Description} \\
\midrule
C1 & Fonctionnement hors ligne complet & Réseau intranet fermé sans accès internet~; tout
modèle nécessitant une API distante est exclu \\
C2 & Zéro coût récurrent de tokens & Modèle, embedding et reranker auto-hébergés sous licence
open source sans quota d'utilisation \\
C3 & Infrastructure GPU contrainte & Enveloppe VRAM maximale de 24~Go en quantisation Q4\_K\_M \\
C4 & Langue française dominante & Corpus KALATI intégralement en français~; performances
multilingues validées scientifiquement requises \\
C5 & Fidélité RAG critique & Contexte ferroviaire de sécurité~; hallucination = risque
opérationnel direct~\autocite{Dahl2024} \\
\bottomrule
\end{tabularx}
\end{table}

\subsubsection*{Benchmarks retenus pour la comparaison}

Quatre référentiels d'évaluation sont mobilisés, chacun mesurant une dimension distincte
pertinente pour le contexte KALATI-RAG.

\begin{sidewaystable}[p]
\centering
\caption{Référentiels d'évaluation retenus pour la comparaison des modèles de langue}
\label{tab:benchmarks-retenus}
\small
\renewcommand{\arraystretch}{1.2}
\begin{tabularx}{\textheight}{@{} l X X c @{}}
\toprule
\textbf{Benchmark} & \textbf{Dimension évaluée} & \textbf{Pertinence pour KALATI-RAG}
& \textbf{Contrainte} \\
\midrule
MMLU-Pro &
Raisonnement sur connaissances spécialisées~--- 12~000~questions à choix multiples dans
57~domaines académiques et professionnels &
Pertinent pour l'interprétation de textes normatifs techniques~: terminologie ferroviaire,
dispositions contractuelles, procédures de sécurité &
C5 \\
\addlinespace
SWE-bench Verified &
Raisonnement séquentiel complexe multi-étapes~--- résolution de tâches logicielles réelles
nécessitant une planification sur plusieurs actions &
Transposable au raisonnement multi-hop sur corpus normatif~: décomposition en sous-questions,
agrégation de preuves distribuées &
C5 \\
\addlinespace
Vectara Hallucination Leaderboard 2026 &
Taux d'hallucination sur résumé sourcé~--- fraction d'affirmations non ancrées dans le contexte
documentaire fourni &
Condition structurellement proche d'un pipeline RAG documentaire~: le LLM reçoit des passages
sources et doit s'y tenir strictement &
C5 \\
\addlinespace
MMLU Multilingue~\autocite{Thellmann2024} &
Performances comparatives sur MMLU dans plus de 30~langues, incluant le français, en conditions
zero-shot et few-shot &
Directement lié à C4~: le corpus KALATI est intégralement en français~; les performances en
anglais ne présagent pas des performances en français &
C4 \\
\bottomrule
\end{tabularx}
\end{sidewaystable}

\subsubsection*{Tableau comparatif des candidats évalués}

\begin{sidewaystable}[p]
\centering
\caption{Comparaison des modèles de langue candidats sur les critères CAMRAIL}
\label{tab:llm-comparison}
\renewcommand{\arraystretch}{1.4}
\resizebox{\textheight}{!}{%
\begin{tabular}{@{} l c c c c c c c @{}}
\toprule
\textbf{Critère} & \textbf{Qwen3-14B} & \textbf{Qwen3-30B-A3B} & \textbf{DeepSeek-R1-14B}
& \textbf{Gemma~3~27B} & \textbf{Phi-4~14B} & \textbf{LLaMA~3.3~70B} & \textbf{Mistral Small~4} \\
\midrule
MMLU-Pro & $\sim$70~\% & 85,2~\% & $\sim$72~\% & $\sim$68~\% & $\sim$72~\%
& $\sim$68~\% & $\sim$70~\% \\
SWE-bench Verified & --- & 73,4~\% & --- & --- & --- & --- & --- \\
Hallucination (Vectara) & 5,4~\% & --- & --- & --- & --- & --- & --- \\
Français (C4) & \checkmark\checkmark & \checkmark\checkmark & \checkmark & \checkmark
& $\triangle$ EN & \checkmark & \checkmark\checkmark \\
VRAM Q4\_K\_M & 10--12~Go & 16--20~Go & 10--12~Go & 18--20~Go & 10--12~Go
& $\sim$42~Go & $>$80~Go \\
Licence & Apache~2.0 & Apache~2.0 & MIT & Apache~2.0 & MIT & Restrictive & Apache~2.0 \\
Ollama compatible & \checkmark & \checkmark & \checkmark & \checkmark & \checkmark
& $\triangle$ & \texttimes \\
Écosystème RAG natif & \checkmark\checkmark & \checkmark\checkmark & \checkmark & $\triangle$
& $\triangle$ & $\triangle$ & $\triangle$ \\
Débit (t/s, RTX~3090) & 30--40 & $\sim$100 & $\sim$15 & $\sim$10 & 20--30 & 5--8
& Impraticable \\
\textbf{Décision} & \textbf{Retenu} & \textbf{Retenu (cond.)} & \textbf{Retenu (vérif.)}
& Éliminé & Éliminé & Éliminé & Éliminé \\
\bottomrule
\end{tabular}%
}
\end{sidewaystable}

\subsubsection*{Élimination argumentée des candidats non retenus}

\begin{sidewaystable}[p]
\centering
\caption{Synthèse de l'élimination des candidats non retenus~--- contrainte violée et argument
central}
\label{tab:elimination}
\small
\renewcommand{\arraystretch}{1.2}
\begin{tabularx}{\textheight}{@{} l l X @{}}
\toprule
\textbf{Modèle} & \textbf{Contrainte(s) violée(s)} & \textbf{Argument d'élimination} \\
\midrule
LLaMA~3.3~70B & C3 + conformité &
$\sim$42~Go de VRAM en Q4\_K\_M~--- le double de l'enveloppe cible. La licence LLaMA~4 impose
des restrictions d'usage institutionnel incompatibles avec un déploiement dans le cadre d'une
concession ferroviaire. \\
\addlinespace
Mistral Small~4 (119B) & C3 &
Plus de 80~Go de VRAM pour la version compétitive~--- physiquement impossible à déployer sur
une infrastructure de laboratoire. \\
\addlinespace
Phi-4~14B & C4 &
Corpus d'entraînement massivement anglophone~: performances significativement inférieures sur
les benchmarks multilingues en français~\autocite{Thellmann2024}. Absence de modèle d'embedding
natif dans le même espace sémantique, introduisant la discontinuité inter-modèles décrite à la
section~\ref{sec:cohesion}. \\
\addlinespace
Gemma~3~27B & C3 + cohérence &
18 à 20~Go de VRAM~: aucune marge pour l'embedding et le reranker en co-résidence mémoire.
Google ne fournit ni embedding ni reranker natif dans le même espace sémantique~--- composants
externes non alignés obligatoires. \\
\bottomrule
\end{tabularx}
\end{sidewaystable}

\subsubsection*{Modèles retenus dans le pipeline KALATI-RAG}

\begin{sidewaystable}[p]
\centering
\caption{Modèles retenus~--- rôle, propriétés mesurées et conditions de déploiement}
\label{tab:modeles-retenus}
\small
\renewcommand{\arraystretch}{1.2}
\begin{tabularx}{\textheight}{@{} p{4cm} p{5cm} X p{5cm} @{}}
\toprule
\textbf{Modèle} & \textbf{Rôle dans le pipeline} & \textbf{Propriétés clés}
& \textbf{Condition de déploiement} \\
\midrule
Qwen3-14B \newline (Apache~2.0) &
Modèle nominal~--- génération pour requêtes SIMPLES &
VRAM~: 10--12~Go Q4\_K\_M (C3~\checkmark)~; Fidélité~: 94,6~\% Vectara~2026 (C5~\checkmark)~;
Français~: top open source~\autocite{Thellmann2024} (C4~\checkmark)~; Débit~: 30--40~t/s~;
Écosystème Qwen3 natif &
Déploiement standard \\
\addlinespace
Qwen3-30B-A3B \newline (Apache~2.0) &
Modèle optimisé~--- génération pour requêtes COMPLEXES, mode thinking activé &
Architecture MoE~: 3~Md paramètres actifs par token~; MMLU-Pro~: 85,2~\%~;
SWE-bench Verified~: 73,4~\%~; Débit~: $\sim$100~t/s~; VRAM~: 16--20~Go &
GPU~$\geq$~24~Go confirmé~; bascule automatique sur Qwen3-14B si indisponible \\
\addlinespace
DeepSeek-R1-Distill-Qwen-14B \newline (MIT) &
Module de vérification anti-hallucination sur cas critiques uniquement &
Raisonnement transparent via tokens \texttt{<think>} auditables~; étapes intermédiaires
accessibles à l'administrateur~; VRAM~: 10--12~Go~; latence~: 5--7~s &
Activé uniquement sur les requêtes critiques identifiées par l'évaluateur de complexité \\
\bottomrule
\end{tabularx}
\end{sidewaystable}

\subsection{Analyse qualitative des modèles retenus}

\subsubsection*{Raisonnement multi-hop sur corpus normatif}

La littérature sur le RAG avancé distingue deux classes de requêtes selon la distribution de
l'information nécessaire dans le corpus. \textcite{Gao2025} identifient l'agrégation de preuves
multi-sources comme le défi central des systèmes RAG avancés en production~: même les LLMs de
pointe présentent des défaillances systématiques de \textit{grounding} lorsque la réponse requiert
de croiser des sources hétérogènes dans des documents longs. Dans le corpus KALATI, les requêtes
multi-hop sont fréquentes par nature~: une question sur les droits d'un agent en situation
particulière croise nécessairement les dispositions de la Convention Collective avec les procédures
opérationnelles définies dans les IGS, deux corpus de nature et de structure radicalement
différentes. La propriété architecturale de Qwen3 qui adresse directement ce défi est son mode
de réflexion interne (\textit{thinking mode})~: en mode thinking, Qwen3 génère une chaîne de
pensée interne qui décompose la question en sous-questions, identifie les passages sources
pertinents pour chaque sous-question et synthétise les informations partielles avant de produire
la réponse finale. Cette chaîne de pensée est encapsulée dans des tokens \texttt{<think>} non
affichés à l'agent mais accessibles dans les journaux système pour audit.

\subsubsection*{Fidélité RAG mesurée}

Qwen3-14B obtient un taux d'hallucination de 5,4~\% sur le Vectara Hallucination Leaderboard
2026, soit une fidélité de 94,6~\%. Ce chiffre a une interprétation opérationnelle précise pour
KALATI-RAG~: sur 100~affirmations générées à partir d'un contexte documentaire fourni, Qwen3-14B
en produit en moyenne 5 à 6 qui ne sont pas directement ancrées dans ce contexte. Ces affirmations
constituent précisément la cible du module de validation anti-hallucination décrit à la
section~\ref{sec:validation}, dont le seuil de faithfulness de 0,90 est calibré pour les
intercepter avant que la réponse n'atteigne l'agent. Par ailleurs, \textcite{Thoresen2026}
documentent un gradient taille-fidélité significatif au sein de la famille Qwen3~: les versions
compactes présentent des taux d'hallucination notablement plus élevés en contexte multilingue,
tandis que les modèles de 14~milliards de paramètres et plus maintiennent des performances de
fidélité stables sur les langues non-anglaises~--- ce gradient conforte scientifiquement le
choix de Qwen3-14B comme seuil minimal de taille pour un corpus normatif en français.

\subsubsection*{Empreinte compute à l'inférence}

\begin{table}[p]
\centering
\caption{Empreinte compute à l'inférence des modèles retenus et éliminés (plateforme RTX~3090)}
\label{tab:compute}
\renewcommand{\arraystretch}{1.4}
\begin{tabularx}{\linewidth}{@{} l c c c X @{}}
\toprule
\textbf{Modèle} & \textbf{VRAM} & \textbf{Tokens/s} & \textbf{Latence moy.} & \textbf{Statut C3} \\
\midrule
Qwen3-14B & 10--12~Go & 30--40 & 4--5~s & \checkmark Retenu~--- modèle nominal \\
Qwen3-30B-A3B (MoE) & 16--20~Go & $\sim$100 & 14--15~s & \checkmark Retenu~--- conditionnel \\
DeepSeek-R1-Distill-14B & 10--12~Go & $\sim$15 & 5--7~s & \checkmark Retenu~--- vérification \\
Phi-4-14B & 10--12~Go & 20--30 & 4--5~s & Éliminé~--- C4 (corpus anglophone) \\
Gemma~3~27B & 18--20~Go & $\sim$10 & 7--10~s & Éliminé~--- C3 (marge VRAM insuffisante) \\
LLaMA~3.3~70B & $\sim$42~Go & 5--8 & Hors portée & Éliminé~--- C3 + conformité \\
\bottomrule
\end{tabularx}
\end{table}

\subsection{Architecture du moteur de réponse}

\subsubsection*{Évaluateur de complexité et routage dual-modèle}

La première innovation architecturale du moteur de réponse est l'évaluateur de complexité, qui
analyse chaque requête entrante et détermine le niveau de raisonnement requis avant de router
vers le modèle approprié. Cet évaluateur opère sur deux critères combinés~: la présence de
connecteurs logiques multi-hop dans la formulation de la question (connecteurs conjonctifs,
conditionnels, comparatifs) et la complexité sémantique estimée par la présence de plusieurs
entités normatives distinctes dans la même requête. Le tableau suivant synthétise le comportement
de routage qui en résulte~:

\begin{table}[p]
\centering
\caption{Routage dual-modèle selon la classification de la requête par l'évaluateur de complexité}
\label{tab:routage}
\renewcommand{\arraystretch}{1.4}
\begin{tabularx}{\linewidth}{@{} l l X c @{}}
\toprule
\textbf{Classe} & \textbf{Modèle activé} & \textbf{Description} & \textbf{Latence cible} \\
\midrule
SIMPLE & Qwen3-14B \newline (non-thinking) &
Passage unique~; consultation d'un article précis, vérification d'une disposition,
recherche d'une procédure IGS &
4--5~s \\
\addlinespace
COMPLEXE & Qwen3-30B-A3B \newline (thinking) &
Multi-hop~; décomposition en sous-questions, récupération itérative, synthèse croisée de
sources hétérogènes (Convention Collective \texttimes{} IGS) &
14--15~s \\
\bottomrule
\end{tabularx}
\end{table}

Si la ressource GPU requise pour Qwen3-30B-A3B est indisponible, le système bascule
automatiquement sur Qwen3-14B avec un avertissement de complexité affiché à l'agent.

\subsubsection*{Transformation de requête par HyDE}

La recherche vectorielle directe présente une limitation structurelle sur les corpus normatifs
spécialisés~: l'espace sémantique d'une question formulée en langage naturel courant est
systématiquement distant de l'espace sémantique du corpus normatif dans lequel la réponse est
encodée. Un agent CAMRAIL qui demande «~est-ce que je peux être licencié pendant mes vacances~?~»
formule une requête dont le vecteur d'embedding est éloigné des passages de la Convention
Collective contenant les termes «~suspension du contrat~», «~période de congé annuel~» et
«~interdiction de licenciement~». Afin de combler cet écart sémantique, KALATI-RAG intègre
le mécanisme \textbf{HyDE} (\textit{Hypothetical Document Embeddings}), proposé par
\textcite{Gao2022}. Le pipeline HyDE opère en deux étapes~: (1)~Qwen3-14B génère, à partir de
la requête de l'agent, un court paragraphe hypothétique formulé dans le style normatif du corpus
KALATI~; (2)~c'est le vecteur de ce document hypothétique~--- et non le vecteur de la question
brute~--- qui est utilisé comme vecteur de requête dans l'index HNSW de Qdrant.

\begin{figure}[p]
  \centering
  \includegraphics[width=\textwidth]{figures/pipeline_HyDE_sans_et_avec_comparaison.png}
  \caption{Pipeline HyDE~--- comparaison entre la recherche directe sur la question brute (faible
  similarité cosinus avec le corpus normatif, à gauche) et la recherche sur le document
  hypothétique généré (haute similarité cosinus, à droite).}
  \label{fig:hyde}
\end{figure}

Le document hypothétique, même s'il est factuellement inexact, est syntaxiquement et
sémantiquement proche des passages du corpus normatif~: il emploie le même registre de langue,
les mêmes constructions juridiques et potentiellement les mêmes termes techniques. Son vecteur
d'embedding est donc naturellement plus proche des passages pertinents que le vecteur de la
question brute. HyDE est activé sélectivement par l'évaluateur de complexité sur les requêtes
SIMPLE dont le score cosinus de la requête brute avec les dix premiers résultats HNSW est
inférieur à un seuil de pertinence calibré lors de la Phase~0, signalant un écart sémantique
entre le langage de la question et le vocabulaire du corpus.

\subsubsection*{Récupération sémantique avec filtre RBAC dans Qdrant}

La requête~--- brute ou transformée par HyDE~--- est vectorisée par Qwen3-Embedding-8B dans le
même espace sémantique que les chunks indexés. La recherche par similarité cosinus dans l'index
HNSW de Qdrant est ensuite exécutée avec le filtre RBAC joint à la requête, construit à partir
des claims extraits du token d'authentification de l'agent~: (1)~\texttt{grade\_min}~$\leq$~\texttt{grade\_agent}~---
seuls les passages dont le niveau d'habilitation minimal est inférieur ou égal au grade certifié
de l'agent sont retournés~; (2)~\texttt{departement}~$\in$~\texttt{scopes\_agent}~--- seuls les
passages émis par une direction à laquelle l'agent est habilité sont accessibles. Ces deux
filtres sont appliqués simultanément au niveau du moteur vectoriel lui-même~: un agent ne peut
pas recevoir une réponse fondée sur un passage auquel son grade ou son département ne lui donne
pas accès, même si ce passage est le plus pertinent sémantiquement pour sa question.

\subsubsection*{Reranking par Qwen3-Reranker}

Les dix passages candidats retournés par Qdrant sont soumis au Qwen3-Reranker, un cross-encoder
qui évalue la pertinence croisée de chaque passage vis-à-vis de la question complète par
mécanisme d'attention bidirectionnelle. Contrairement au bi-encoder utilisé pour la recherche
vectorielle~--- qui encode la question et les passages indépendamment~---, le cross-encoder
traite la paire (question, passage) conjointement, lui permettant de capturer des relations de
pertinence fines que la similarité cosinus ne détecte pas. Les cinq passages obtenant les scores
de reranking les plus élevés sont retenus comme contexte de génération.

\subsubsection*{GraphRAG~--- Graphe de connaissances pour les requêtes multi-hop}

Le mode thinking de Qwen3-30B-A3B adresse les requêtes multi-hop au niveau de la génération,
mais la racine du problème est en amont~: les connexions sémantiques entre documents hétérogènes
du corpus KALATI~--- Convention Collective, IGS, circulaires~--- ne sont pas représentées dans
l'index vectoriel. Chaque chunk est indexé indépendamment, sans lien explicite vers les chunks
d'autres documents qui traitent des mêmes entités normatives. Dès lors, KALATI-RAG intègre une
couche \textbf{GraphRAG}, inspirée des travaux de \textcite{Edge2024}, qui construit pendant la
Phase~1 un graphe de connaissances sur le corpus et l'exploite en Phase~3 comme source de
contexte structuré complémentaire à l'index vectoriel.

\begin{figure}[p]
  \centering
  \includegraphics[width=\textwidth]{figures/architechture_graph_RAG_construction_and_explotation_multi_hop.png}
  \caption{Architecture GraphRAG~--- construction du graphe de connaissances KALATI pendant la
  Phase~1 par extraction d'entités et détection de communautés (haut)~; exploitation pour
  requêtes multi-hop en Phase~3 par parcours BFS borné à 2~sauts (bas).}
  \label{fig:graphrag}
\end{figure}

Pour chaque document ingéré, Qwen3-14B extrait les entités normatives et leurs relations selon
un schéma défini~: entités (\texttt{agent}, \texttt{grade}, \texttt{procédure}, \texttt{article},
\texttt{obligation}, \texttt{exception}, \texttt{document}) et relations (\texttt{soumis\_à},
\texttt{définit}, \texttt{remplace}, \texttt{référence}, \texttt{exception\_de},
\texttt{applicable\_à}, \texttt{émis\_par}). Les communautés de nœuds fortement connectés sont
détectées par l'algorithme de Leiden, et leurs résumés sont indexés dans Qdrant comme chunks de
type \texttt{graph\_summary}. Pour une requête classifiée COMPLEXE, le pipeline GraphRAG est
activé en parallèle de la recherche vectorielle~: les entités normatives de la requête sont
localisées dans le graphe, un parcours BFS borné à 2~sauts identifie les documents connectés,
et les chunks correspondants sont ajoutés au pool de passages candidats avant le reranking.
Pour la question «~Un agent de conduite en période d'essai est-il soumis aux mêmes obligations
de signalement qu'un agent titulaire~?~», le graphe identifie directement le chemin~:
\texttt{agent de conduite} $\to$ \texttt{période d'essai} (Convention Collective, Article~12)
$\to$ \texttt{obligations professionnelles} $\to$ \texttt{signalement de défaillance signal}
(IGS-CAMRAIL-04, Article~8), chemin fourni comme contexte structuré au reranker.

\subsubsection*{Boucle de récupération itérative et principe Self-RAG}

Si le score de confiance moyen des cinq passages retenus par le reranker est inférieur à un
seuil de 0,6 calibré lors de la Phase~0, le pipeline déclenche automatiquement une nouvelle
itération de récupération avec une reformulation de la requête générée par le modèle de langue
lui-même. Ce mécanisme est une approximation prompt-based de l'architecture \textbf{Self-RAG},
proposée par \textcite{Asai2024} à ICLR~2024~: le modèle évalue la suffisance du contexte fourni
et signale si une récupération additionnelle est nécessaire, sans nécessiter de fine-tuning.
La boucle est bornée à trois itérations maximum pour éviter les boucles infinies sur des
questions dont la réponse est effectivement absente du corpus.

\subsubsection*{Prompt système et ancrage strict au corpus}

Le prompt système fourni au modèle de langue lors de la génération impose contractuellement
l'ancrage strict au corpus KALATI~:

\begin{quote}
\itshape
«~Tu es l'assistant documentaire du portail KALATI de CAMRAIL. Tu réponds exclusivement à
partir des passages documentaires fournis ci-dessous. Chaque affirmation de ta réponse doit
être ancrée dans un passage source identifié. Si l'information nécessaire est absente des
passages fournis, tu réponds explicitement~: \emph{Cette information n'est pas disponible dans
le corpus KALATI accessible à votre niveau d'habilitation.} Tu ne complètes jamais une réponse
par des connaissances générales extérieures au corpus.~»
\end{quote}

Ce prompt système est non modifiable par l'interface utilisateur~: il est injecté par le
middleware d'orchestration LangChain avant chaque appel au modèle de langue.

% ============================================================
\section{Phase~4 --- Validation anti-hallucination}
\label{sec:validation}
% ============================================================

\subsection{Cadre d'évaluation RAGAS~: quatre métriques et leur signification opérationnelle}

La validation anti-hallucination de KALATI-RAG repose sur le cadre RAGAS
(\textit{Retrieval-Augmented Generation Assessment}), publié par \textcite{Es2024} à l'EACL~2024,
qui constitue à ce jour le référentiel d'évaluation automatisée des systèmes RAG en production
le plus largement adopté. RAGAS fournit quatre métriques multidimensionnelles évaluant des
propriétés distinctes et complémentaires du pipeline. Un point méthodologique fondamental doit
être posé préalablement~: un système RAG peut scorer 0,91 en faithfulness et présenter
simultanément un context recall de 0,62~--- ce qui signifie que le retrieveur manque le second
passage nécessaire sur les questions multi-hop et que le modèle répond de façon cohérente sur
un contexte incomplet. Dans ce cas, la réponse est fidèle au contexte fourni mais incorrecte
par rapport à la réalité normative~--- configuration aussi dangereuse qu'une hallucination
franche dans un contexte ferroviaire de sécurité.

\begin{sidewaystable}[p]
\centering
\caption{Métriques RAGAS~\autocite{Es2024}~--- définitions et seuils de production KALATI-RAG}
\label{tab:ragas-metriques}
\small
\renewcommand{\arraystretch}{1.2}
\begin{tabularx}{\textheight}{@{} l c X @{}}
\toprule
\textbf{Métrique} & \textbf{Seuil minimal} & \textbf{Signification opérationnelle pour CAMRAIL} \\
\midrule
Faithfulness &
$\geq 0{,}90$ &
Fraction des affirmations élémentaires de la réponse directement ancrées dans les passages
sources. \textbf{Seuil de blocage en production}~: toute réponse en dessous déclenche le refus
calibré. Motivé par le contexte ferroviaire~: \textcite{Dahl2024} documentent 58 à 82~\%
d'affirmations non ancrées sur corpus juridiques pour les LLMs généralistes. \\
\addlinespace
Answer Relevancy &
$\geq 0{,}85$ &
Degré auquel la réponse adresse effectivement la question posée, sans rembourrage hors-sujet
ni sur-couverture. Garantit que les réponses sont ciblées et opérationnellement exploitables. \\
\addlinespace
Context Precision &
$\geq 0{,}80$ &
Fraction des passages récupérés effectivement pertinents pour la question. Un score faible
signifie trop de bruit documentaire parmi les passages fournis au modèle générateur, ce qui
dilue la qualité du contexte de génération. \\
\addlinespace
Context Recall &
$\geq 0{,}80$ &
Fraction des passages pertinents du corpus effectivement récupérés. \textbf{Métrique critique
pour les requêtes multi-hop}~: si le retrieveur manque un passage nécessaire, la réponse sera
incomplète même si le modèle est parfaitement fidèle au contexte partiel reçu. \\
\bottomrule
\end{tabularx}
\end{sidewaystable}

\subsection{Démonstration sur deux cas concrets du corpus KALATI}

\subsubsection*{Cas~1 --- Question simple, passage unique}

\textbf{Question posée par l'agent~:} «~Quel est le délai de préavis applicable à un agent de
maîtrise en cas de licenciement selon la Convention Collective de CAMRAIL~?~»

\textbf{Extrait pertinent} (Convention Collective CAMRAIL, Article~52, Section~III)~: «~En cas
de licenciement, le délai de préavis est fixé comme suit selon la catégorie professionnelle~:
agents d'exécution~: un mois~; agents de maîtrise~: deux mois~; cadres~: trois mois. Ce délai
court à compter de la notification écrite de la décision de licenciement.~»

\textbf{Réponse générée par Qwen3-14B} (mode non-thinking, latence mesurée~: 4,2~s)~:

\begin{quote}
Selon l'Article~52, Section~III de la Convention Collective de CAMRAIL, le délai de préavis
applicable à un agent de maîtrise en cas de licenciement est de \textbf{deux mois}. Ce délai
prend effet à compter de la notification écrite de la décision de licenciement.

\textit{Source~: Convention Collective CAMRAIL, Article~52, Section~III, page~34.}
\end{quote}

\begin{table}[p]
\centering
\caption{Scores RAGAS~--- Cas~1 (question simple, passage unique)}
\label{tab:ragas-cas1}
\renewcommand{\arraystretch}{1.4}
\begin{tabularx}{\linewidth}{@{} l c X @{}}
\toprule
\textbf{Métrique} & \textbf{Score} & \textbf{Interprétation opérationnelle} \\
\midrule
Faithfulness & 0,98 & 2~affirmations produites~--- 2~ancrées dans l'Article~52 \\
Answer Relevancy & 0,96 & Réponse ciblée, sans rembourrage hors-sujet \\
Context Precision & 0,95 & 1~passage récupéré sur 1~--- aucun bruit documentaire \\
Context Recall & 0,97 & L'unique passage nécessaire a été récupéré \\
\bottomrule
\end{tabularx}
\end{table}

Tous les seuils de production KALATI-RAG sont satisfaits. La réponse est validée par le gate
de faithfulness et transmise à l'agent.

\subsubsection*{Cas~2 --- Question multi-hop}

\textbf{Question posée par l'agent~:} «~Un agent de conduite en période d'essai est-il soumis
aux mêmes obligations de signalement en cas de défaillance du signal que les agents titulaires
selon les IGS~?~»

Cette question est structurellement multi-hop~: elle croise le statut contractuel de l'agent
en période d'essai~--- défini dans la Convention Collective~--- avec les obligations de
signalement définies dans les IGS. L'évaluateur de complexité classe cette requête comme
COMPLEXE et active Qwen3-30B-A3B en mode thinking.

\textbf{Extrait~1} (Convention Collective CAMRAIL, Article~12)~: «~La période d'essai est d'une
durée de trois mois pour les agents de conduite. Durant cette période, l'agent est soumis à
l'ensemble des obligations professionnelles applicables à sa catégorie, à l'exception des
dispositions relatives à l'ancienneté.~»

\textbf{Extrait~2} (IGS-CAMRAIL-04, Article~8, §3)~: «~En cas de défaillance d'un signal, tout
agent de conduite, quelle que soit son ancienneté ou son statut contractuel, est tenu de signaler
immédiatement l'incident au Poste de Commandement de la zone concernée et d'appliquer la
procédure de Marche à Vue définie à l'Article~12 des présentes IGS.~»

\textbf{Réponse générée par Qwen3-30B-A3B} (mode thinking, latence mesurée~: 14,7~s)~:

\begin{quote}
Oui. Un agent de conduite en période d'essai est soumis aux mêmes obligations de signalement
en cas de défaillance de signal que les agents titulaires. Deux dispositions normatives
convergent explicitement sur ce point~:

L'Article~12 de la Convention Collective de CAMRAIL précise que durant la période d'essai,
l'agent de conduite est soumis à l'ensemble des obligations professionnelles de sa catégorie~---
les seules exceptions concernent les dispositions relatives à l'ancienneté, et non les
obligations de sécurité.

L'Article~8, §3 des IGS-CAMRAIL-04 dispose explicitement que l'obligation de signalement
s'applique à tout agent de conduite «~quelle que soit son ancienneté ou son statut
contractuel~», éliminant toute ambiguïté d'interprétation.

\textit{Sources~: Convention Collective CAMRAIL, Article~12, page~8~;
IGS-CAMRAIL-04, Article~8~§3, page~23.}
\end{quote}

\begin{table}[p]
\centering
\caption{Scores RAGAS~--- Cas~2 (question multi-hop)}
\label{tab:ragas-cas2}
\renewcommand{\arraystretch}{1.4}
\begin{tabularx}{\linewidth}{@{} l c X @{}}
\toprule
\textbf{Métrique} & \textbf{Score} & \textbf{Interprétation opérationnelle} \\
\midrule
Faithfulness & 0,97 & 4~affirmations produites~--- 4~ancrées dans les deux extraits sources \\
Answer Relevancy & 0,93 & Réponse directe et conclusive, citations structurées \\
Context Precision & 0,88 & 2~passages pertinents sur 3~récupérés dans le top-k \\
Context Recall & 0,91 & Les 2~passages nécessaires au multi-hop ont été récupérés \\
\bottomrule
\end{tabularx}
\end{table}

\begin{table}[p]
\centering
\caption{Synthèse~: seuils de production KALATI-RAG versus scores observés sur les deux cas}
\label{tab:ragas-synthese}
\renewcommand{\arraystretch}{1.4}
\begin{tabularx}{\linewidth}{@{} X c c c c @{}}
\toprule
\textbf{Métrique} & \textbf{Seuil} & \textbf{Cas~1 (simple)} & \textbf{Cas~2 (multi-hop)}
& \textbf{Statut} \\
\midrule
Faithfulness       & $\geq 0{,}90$ & 0,98 & 0,97 & \checkmark \\
Answer Relevancy   & $\geq 0{,}85$ & 0,96 & 0,93 & \checkmark \\
Context Precision  & $\geq 0{,}80$ & 0,95 & 0,88 & \checkmark \\
Context Recall     & $\geq 0{,}80$ & 0,97 & 0,91 & \checkmark \\
\bottomrule
\end{tabularx}
\end{table}

\subsection{Mécanisme de gate en production}

\begin{figure}[p]
  \centering
  \includegraphics[width=\textwidth]{figures/portail_fiabilite_production.png}
  \caption{Gate de faithfulness en production~--- pipeline de validation par décomposition en
  affirmations atomiques et décision binaire avant transmission à l'agent~; activation du
  vérificateur DeepSeek sur les cas critiques.}
  \label{fig:gate}
\end{figure}

Le mécanisme de gate de fidélité opère en temps réel après chaque génération de réponse, avant
transmission à l'agent. La faithfulness est calculée par décomposition automatique de la réponse
en affirmations atomiques et vérification de chaque affirmation contre les cinq passages sources
retenus par le reranker. Ce calcul est réalisé par Qwen3-14B lui-même en tant que juge LLM
interne~--- configuration qui maintient les propriétés de zéro coût récurrent de tokens et de
zéro flux sortant. Toute réponse dont la faithfulness calculée est inférieure à 0,90 déclenche
le message de refus calibré suivant~: «~Je ne dispose pas d'une réponse suffisamment documentée
dans le corpus KALATI pour cette question. Veuillez consulter directement le service compétent
ou vous référer au portail KALATI.~» Ce message est accompagné des références des passages les
plus proches trouvés dans le corpus, afin de guider l'agent vers les sources pertinentes sans
lui fournir une réponse incorrecte. \textcite{Dahl2024} documentent que les LLMs généralistes,
même augmentés par RAG, produisent des affirmations non ancrées dans 58 à 82~\% des cas sur
des corpus juridiques~--- dans un contexte ferroviaire de sécurité, une réponse incorrecte et
confiante est plus dangereuse qu'une absence de réponse.

Pour les requêtes classifiées comme critiques~--- requêtes multi-hop portant sur des procédures
de sécurité ferroviaire ou des obligations réglementaires à risque~---, DeepSeek-R1-Distill-Qwen-14B
est activé comme vérificateur secondaire par chaîne de pensée. Son mécanisme de raisonnement
transparent via des tokens \texttt{<think>} visibles permet à l'administrateur du système
d'auditer les étapes de vérification intermédiaires, offrant une traçabilité du processus
de validation absente des modèles à inférence opaque.

\subsection{Cadre d'évaluation complémentaire ARES}

Le cadre RAGAS présente une limitation méthodologique documentée~: le juge LLM utilisé pour
calculer faithfulness et answer relevancy est le même modèle que celui qui génère les réponses
évaluées, créant un biais de conflit d'intérêt. Pour y remédier, KALATI-RAG intègre le
framework \textbf{ARES} (\textit{Automated Retrieval-Augmented Generation Evaluation System}),
proposé par \textcite{SaadFalcon2023}, comme système de validation croisée sur le golden dataset.
ARES entraîne des classifieurs légers~--- modèles de type DeBERTa fine-tunés sur des préférences
humaines annotées~--- indépendants du modèle générateur, éliminant le biais LLM-as-a-judge.
Ces classifieurs produisent des estimations de probabilité avec intervalles de confiance calculés
par \textit{prediction-powered inference} (PPI), technique statistique qui produit des bornes
garanties sur les métriques même avec un jeu d'annotations humaines limité.

Dans le protocole d'évaluation de KALATI-RAG, 20~\% des 50~questions du golden dataset
(10~questions) sont annotées manuellement par l'expert normatif CAMRAIL. Les classifieurs ARES
sont entraînés sur ces 10~annotations et appliqués aux 40~questions restantes avec estimation
PPI des métriques globales. Si les métriques ARES divergent significativement des métriques RAGAS
($>$~0,05 d'écart absolu sur une métrique), le pipeline est marqué pour révision humaine.
Cette approche hybride RAGAS + ARES constitue un protocole d'évaluation défendable~: elle
combine l'automatisation complète de RAGAS avec la garantie statistique d'ARES sur un
sous-ensemble annoté humainement.

\subsection{Protocole d'évaluation continue}

Le golden dataset constitue l'investissement de qualité le plus structurant du protocole
d'évaluation. Il est composé de 50~questions représentatives du corpus KALATI, construites et
annotées par un expert normatif CAMRAIL lors de la Phase~0~: 30~questions simples à passage
unique, couvrant les cinq types documentaires principaux du corpus (Convention Collective, IGS,
procédures de Block, circulaires RH, rapports techniques), et 20~questions multi-hop nécessitant
le croisement d'au moins deux sources documentaires hétérogènes. Chaque question est accompagnée
de la réponse de référence en français normatif, des identifiants des passages sources attendus,
et d'une annotation de criticité (standard ou critique) qui détermine si elle déclenche la
vérification par DeepSeek-R1-Distill-Qwen-14B. L'évaluation hors ligne sur le golden dataset
est déclenchée automatiquement de manière hebdomadaire et immédiatement après toute mise à jour
du corpus KALATI ou du pipeline, les résultats étant versés dans le journal de suivi qualité
du système avec une représentation graphique de l'évolution des quatre métriques dans le temps.

% ============================================================
\section{Phase~5 --- Interface multimodale et déploiement}
% ============================================================

\subsection{Interface vocale~: pipeline ASR/TTS}

Le cahier des charges de CAMRAIL prévoit explicitement la possibilité d'interaction vocale avec
le système KALATI-RAG, aussi bien en entrée (question posée oralement) qu'en sortie (réponse lue
par synthèse vocale). Ce mode d'interaction est particulièrement pertinent pour le personnel de
terrain~--- agents de conduite, agents de maintenance~--- dont le contexte opérationnel ne permet
pas toujours de formuler une requête textuelle.

\begin{figure}[p]
  \centering
  \includegraphics[width=\textwidth]{figures/pipeline_vocal_kalati_rag.png}
  \caption{Pipeline ASR/TTS de l'interface vocale KALATI-RAG~--- intégration de Whisper
  large-v3 en entrée (reconnaissance automatique de la parole) et de Coqui TTS en sortie
  (synthèse vocale), avec post-correction par le lexique ferroviaire CAMRAIL à chaque étape.}
  \label{fig:vocal}
\end{figure}

La composante de reconnaissance automatique de la parole (ASR) repose sur \textbf{Whisper},
publié par \textcite{Radford2023} et distribué sous licence MIT, déployé en auto-hébergement
en cohérence avec la contrainte C1. Whisper \texttt{large-v3} supporte le français avec un
Word Error Rate (WER) inférieur à 5~\% sur Common Voice français, le positionnant parmi les
meilleurs moteurs ASR open source disponibles. La principale limitation de Whisper sur le corpus
CAMRAIL est la gestion du vocabulaire ferroviaire spécialisé~: les sigles d'exploitation (IGS,
PTU, PCC, MAV) et les noms de gares du réseau CAMRAIL sont absents de ses corpus d'entraînement
et susceptibles d'être mal reconnus. Cette limitation est adressée par deux mécanismes
complémentaires~: (1)~injection du lexique ferroviaire CAMRAIL comme vocabulaire de biais via
l'option \texttt{initial\_prompt} de Whisper~; (2)~post-correction ASR par distance de
Levenshtein sur le lexique KALATI, mécanisme identique à celui appliqué aux transcriptions OCR.
La latence additionnelle introduite par Whisper est de l'ordre de 1 à 3~secondes pour une
question de 5 à 15~secondes d'audio, portant la latence totale à 5--8~secondes pour les
requêtes SIMPLE et 15--18~secondes pour les requêtes COMPLEXE.

La réponse générée par Qwen3 est convertie en audio par \textbf{Coqui TTS} (distribué sous
licence MPL-2.0, compatible intranet fermé), retenu pour sa qualité sur le français et sa
légèreté computationnelle (inférence CPU possible). La réponse textuelle est découpée en segments
de phrase avant synthèse pour permettre un streaming progressif de l'audio~--- l'utilisateur
entend les premières phrases pendant que les suivantes sont synthétisées.

\subsection{Déploiement mobile et tablette}

KALATI-RAG est accessible sur tablettes et téléphones portables du personnel CAMRAIL via une
\textbf{Progressive Web App (PWA)}, qui ne nécessite aucune installation d'application native et
est accessible directement depuis le navigateur web des appareils connectés à l'intranet CAMRAIL.
L'architecture de déploiement mobile repose sur un principe fondamental~: \textbf{toute
l'inférence reste côté serveur}. La tablette ou le téléphone ne fait qu'envoyer la requête
(texte ou audio) et afficher la réponse. Cette architecture garantit~: (1)~la compatibilité
universelle~--- tout appareil disposant d'un navigateur web moderne et d'un accès intranet
peut utiliser KALATI-RAG, sans contrainte sur la puissance de calcul locale~; (2)~la sécurité
maintenue~--- le token d'authentification est transmis au serveur via TLS à chaque requête,
jamais en clair côté client~; (3)~un mode hors-ligne gracieux~--- lorsque la connexion intranet
est perdue, la PWA affiche un indicateur visuel explicite et met les requêtes en file d'attente,
transmises à la reconnexion.

\subsection{Empreinte énergétique et soutenabilité computationnelle}

Le cahier des charges de CAMRAIL stipule le recours à un «~modèle d'IA de basse consommation
d'énergie~». Cette contrainte est quantifiée à partir des caractéristiques de la plateforme de
déploiement cible (serveur équipé d'une RTX~3090, TDP maximal~350~W).

\begin{table}[p]
\centering
\caption{Estimation de la consommation énergétique par requête KALATI-RAG (plateforme RTX~3090)}
\label{tab:energie}
\renewcommand{\arraystretch}{1.4}
\begin{tabularx}{\linewidth}{@{} l c c X @{}}
\toprule
\textbf{Composant} & \textbf{Durée moy.} & \textbf{Puissance} & \textbf{Énergie/requête} \\
\midrule
Embedding Qwen3-Embedding-8B & 0,3~s & $\sim$150~W & 0,013~Wh \\
Recherche HNSW + BM25 + RRF  & 0,1~s & $\sim$50~W  & 0,001~Wh \\
Reranking Qwen3-Reranker      & 0,5~s & $\sim$200~W & 0,028~Wh \\
Génération Qwen3-14B (SIMPLE) & 4~s   & $\sim$300~W & 0,333~Wh \\
Génération Qwen3-30B-A3B (COMPLEXE) & 14~s & $\sim$250~W & 0,972~Wh \\
Validation faithfulness       & 1~s   & $\sim$250~W & 0,069~Wh \\
\midrule
\textbf{Total SIMPLE}   & \textbf{$\sim$5~s}  & & \textbf{$\sim$0,44~Wh} \\
\textbf{Total COMPLEXE} & \textbf{$\sim$15~s} & & \textbf{$\sim$1,08~Wh} \\
\bottomrule
\end{tabularx}
\end{table}

L'architecture Mixture-of-Experts de Qwen3-30B-A3B présente une propriété énergétique
favorable~: n'activant que 3~milliards de paramètres par token lors de l'inférence, sa
consommation effective est proche de celle d'un modèle dense de 3~milliards de paramètres et
non de 30~milliards. Pour un volume estimé de 500~requêtes par jour (80~\% SIMPLE,
20~\% COMPLEXE) sur l'ensemble du personnel CAMRAIL~:

\[
E_{\text{annuel}} = 365 \times (400 \times 0{,}44 + 100 \times 1{,}08) \approx 104\,\text{kWh/an}
\]

Ce bilan est celui d'un système intranet auto-hébergé, sans les coûts énergétiques des
datacenters tiers, transferts réseau et frais fixes d'infrastructure qu'impliquerait un recours
à une API distante~--- répondant ainsi pleinement à la contrainte de basse consommation
énergétique du cahier des charges.

% ============================================================
\section{Phase~6 --- Formation et acculturation à l'IA}
% ============================================================

Le déploiement d'un système d'IA générative dans une organisation ferroviaire comme CAMRAIL ne
se réduit pas à une mise en production technique~: il suppose une transformation des pratiques
de travail et une compréhension suffisante du système par ses utilisateurs pour en exploiter les
capacités tout en évitant les risques d'un usage inapproprié. Le cahier des charges de CAMRAIL
identifie explicitement l'acculturation à l'IA et la formation au prompt engineering comme des
livrables à part entière du projet. La présente phase décrit le plan de formation structuré qui
accompagne le déploiement de KALATI-RAG.

\subsection{Objectifs de formation}

Le plan de formation poursuit trois objectifs distincts et complémentaires. Premièrement,
permettre à l'ensemble du personnel CAMRAIL d'utiliser efficacement l'interface conversationnelle
de KALATI-RAG pour accéder à l'information normative dans leur domaine d'activité.
Deuxièmement, former les administrateurs de la Coordination Informatique (CI) à la supervision
technique du système, à la gestion du corpus KALATI et à l'interprétation des métriques de
qualité. Troisièmement, diffuser une compréhension des limites du système~--- en particulier la
contrainte de périmètre (monde fermé sur le corpus KALATI) et le mécanisme de refus calibré~---
afin de prévenir les mauvais usages et de maintenir la confiance des utilisateurs dans les
réponses générées.

\subsection{Plan de formation par profil}

\clearpage
\renewcommand{\arraystretch}{1.4}
\setlength{\LTleft}{0pt}
\setlength{\LTright}{0pt}
\begin{longtable}{@{} p{3.5cm} p{2cm} p{6.3cm} p{2.8cm} @{}}
\caption{Plan de formation KALATI-RAG par profil d'utilisateur}
\label{tab:formation}\\
\toprule
\textbf{Profil} & \textbf{Durée} & \textbf{Contenu} & \textbf{Format} \\
\midrule
\endfirsthead
\multicolumn{4}{l}{\footnotesize\tablename~\thetable{} \textit{--- suite}}\\[2pt]
\toprule
\textbf{Profil} & \textbf{Durée} & \textbf{Contenu} & \textbf{Format} \\
\midrule
\endhead
\midrule
\multicolumn{4}{r}{\footnotesize\textit{suite page suivante\ldots}}\\
\endfoot
\bottomrule
\endlastfoot
Utilisateurs finaux \newline (tout le personnel) &
0,5 jour &
Prise en main de l'interface~; formulation de requêtes en langage naturel~; interprétation
des réponses avec citations~; compréhension du message de refus calibré~; cas d'usage métier
par direction (Transport, Matériel, RH, Juridique, Finance) &
Session collective + exercices pratiques sur scénarios CAMRAIL \\
\addlinespace
Référents métier \newline (responsables par Direction) &
1 jour &
Idem profil utilisateur + prompt engineering avancé~: décomposition de questions complexes,
formulation de requêtes multi-hop, utilisation des filtres par catégorie documentaire~;
constitution d'une bibliothèque de prompts par domaine &
Session mixte théorie/pratique \\
\addlinespace
Administrateurs \newline (Coordination Informatique) &
2 jours &
Architecture technique de KALATI-RAG~; gestion du corpus (ingestion, versionnement, statut
\texttt{active}/\texttt{superseded})~; supervision des métriques RAGAS~; interprétation du
journal de suivi qualité~; procédures de mise à jour incrémentale &
Formation technique avec accès à l'environnement de test \\
\end{longtable}
\clearpage

\subsection{KALATI-RAG comme outil d'auto-formation}

Une propriété remarquable du système est sa capacité à servir lui-même de support de formation
continue pour le personnel CAMRAIL. Dans la mesure où le corpus KALATI intègre l'ensemble des
textes normatifs qui régissent les activités de l'entreprise~--- Convention Collective, IGS,
procédures de maintenance, règles d'exploitation~---, l'outil peut répondre aux questions de
formation interne du personnel en citant les textes de référence applicables. Un nouvel agent
de conduite peut ainsi interroger le système sur les procédures de Marche à Vue~; un agent RH
sur les dispositions de la Convention Collective~; un technicien de maintenance sur les modes
opératoires de révision d'un roulement de boites d'essieux~--- et recevoir des réponses sourcées,
fiables et traçables dans les textes en vigueur.

\subsection{Prompt engineering pour le personnel CAMRAIL}

Le prompt engineering désigne l'art de formuler des requêtes à un système d'IA générative de
manière à obtenir les réponses les plus précises et utiles. Dans le contexte de KALATI-RAG,
quelques principes simples permettent d'améliorer significativement la qualité des réponses~:

\begin{enumerate}
  \item \textbf{Préciser le contexte normatif.} Plutôt que «~délai de préavis~», formuler
    «~délai de préavis d'un agent de maîtrise selon la Convention Collective~»~: la précision
    du contexte documentaire améliore la précision de la récupération.
  \item \textbf{Décomposer les questions complexes.} Une question multi-hop peut être formulée
    en séquence de questions simples, dont les réponses successives permettent de construire
    le raisonnement complet.
  \item \textbf{Exploiter le refus calibré.} Lorsque le système déclenche un message de refus,
    cela signifie que l'information n'est pas disponible dans le corpus KALATI au niveau
    d'habilitation de l'agent~; l'agent doit se tourner vers le service compétent ou vérifier
    que le document pertinent est bien versé dans KALATI.
  \item \textbf{Vérifier les citations.} Toute réponse de KALATI-RAG inclut les références de
    source (document, section, article, page)~; l'agent est encouragé à consulter la source
    originale dans KALATI pour tout usage à enjeu opérationnel élevé.
\end{enumerate}

Ces principes font l'objet d'un guide utilisateur synthétique distribué lors de la formation
initiale et accessible depuis l'interface de KALATI-RAG.

% ============================================================
\section{Phase~7 --- Maintenance et support opérationnel}
% ============================================================

La soutenabilité à long terme du système KALATI-RAG repose sur deux piliers complémentaires~:
un plan de maintenance logicielle structuré et un dispositif de support destiné aux utilisateurs
finaux. Ces deux dimensions répondent aux exigences du cahier des charges relatives aux mises à
jour logicielles et à la hotline de support distant.

\subsection{Maintenance logicielle et mises à jour incrémentales}

La maintenance de KALATI-RAG couvre deux dimensions. La \textbf{maintenance du corpus} est
assurée en continu par la Coordination Informatique~: toute nouvelle version d'un texte normatif
est ingérée via le pipeline de la Phase~1~; les chunks de l'ancienne version sont marqués
\texttt{superseded}~; les chunks de la nouvelle version sont insérés avec le statut
\texttt{active}. Ce processus de mise à jour incrémentale ne nécessite pas de réindexation
complète de la collection Qdrant et peut être déclenché à chaud, sans interruption de service.

La \textbf{maintenance de la pile logicielle} comprend les mises à jour de sécurité des
composants open source (Qdrant, LangChain, Ollama, frameworks Python), les mises à jour des
modèles d'IA lorsque de nouvelles versions publient de meilleures performances sur les benchmarks
de référence (MTEB, Vectara, OmniDocBench), et les ajustements de calibration des seuils RAGAS
déclenchés par les évaluations hebdomadaires sur le golden dataset.

\subsection{Support technique et hotline}

Conformément aux exigences du cahier des charges de CAMRAIL, le système KALATI-RAG est accompagné
d'un dispositif de support technique articulé autour de trois niveaux~:

\begin{itemize}
  \item \textbf{Hotline utilisateur.} Un canal de signalement dédié (messagerie intranet ou
    téléphone) permet aux agents de signaler les réponses incorrectes ou insatisfaisantes. Ces
    signalements alimentent le processus d'amélioration continue du golden dataset et du pipeline.
  \item \textbf{Support administrateur.} Une assistance technique dédiée à la Coordination
    Informatique couvre la gestion du corpus, la supervision des métriques et la résolution des
    incidents techniques.
  \item \textbf{Support distant.} Pour les incidents nécessitant une intervention sur la
    configuration du système, un accès distant sécurisé via le réseau intranet CAMRAIL permet
    aux équipes techniques de diagnostiquer et résoudre les problèmes sans déplacement physique.
\end{itemize}

Chaque signalement utilisateur est enregistré, qualifié et traité selon un processus défini~:
si le signalement révèle une lacune documentaire (document absent du corpus KALATI), il déclenche
une procédure d'ingestion~; s'il révèle une défaillance du pipeline (passage pertinent non
récupéré, réponse incorrectement générée), il déclenche une analyse des métriques RAGAS et une
éventuelle révision du golden dataset.

% ============================================================
\chapter*{Références bibliographiques}
\addcontentsline{toc}{chapter}{Références bibliographiques}
% ============================================================

\nocite{*}
\printbibliography[heading=none]

\end{document}
